\documentclass{article}

\begin{document}

\title{GPU Programming with Directives Exercise \
  Logging on, Compiling and Running}

\maketitle

\section{Introduction}


The purpose of this exercise is to check you can compile and run
simple GPU programs on ARCHER2.

\section{Connecting to ARCHER2}

Log on to ARCHER2: \verb+ssh username@login.archer2.ac.uk+ \\

Note: for this exercise, it is important to use the work file system,
i.e. \verb+cd /work/project/project/username+, and not the home file system.

\section{Compiling and Running}
The code is contained in saxyp.tar. The tar file can be fetched from GitHub by
cloning the course repository with the following commands:

\begin{verbatim}
git clone https://github.com/EPCCed/archer2-GPU-directives.git
cd archer2-GPU-directives
git checkout 2025-04-22
cd Exercises
\end{verbatim}

Alternatively, the file can be found on ARCHER2 and copied into your
/work/ directory with the command:

\begin{verbatim}
cp /work/z19/shared/GPUdir/saxpy.tar .
\end{verbatim}

Unpack the file with the command: \verb+tar -xvf saxpy.tar+

\section{Compiling and Running}

Since this exercise will involve offloading to GPUs with OpenMP
directives, certain modules must be loaded prior to compiling the
code:

\begin{verbatim}
module load PrgEnv-amd
module load rocm
module load craype-accel-amd-gfx90a
module load craype-x86-milan
\end{verbatim}

There are two examples, one using HIP and the other OpenMP. To compile and run the HIP example:

\begin{verbatim}
cd saxpy
make -f Makefile-hip
sbatch hipsaxpy.slurm
\end{verbatim}

and similarly for the OpenMP example.

\end{document}
